\section{Методология исследования}

\subsection{Подготовка признаков}

Целевая переменная была упорядочена от класса \textit{Hopeless} до класса
\textit{Prospering}. Каждому классу сопоставлялся числовой ранг от 0 до 9. Это
позволило использовать ранговые корреляции при отборе состояний, но итоговая
задача оставалась многоклассовой классификацией.

Причины предварительно проверялись на мультиколлинеарность. Использовался
итеративный VIF-фильтр: пока максимальный VIF был не меньше 8, удалялся признак
с наибольшим VIF. В результате были исключены два общественных признака:
\textit{издержки сообщества на окружающую среду} и \textit{волатильность
потребительских цен в сообществе}. Осталось 8 причин.

\subsection{Отбор состояний}

Для каждого состояния рассчитывалась корреляция Спирмена с рангом счастья.
Кандидатами считались состояния с достаточно сильной монотонной связью:
\begin{equation}
    |\rho_S(S_j, rank(y))| \ge 0{,}40.
\end{equation}
После этого состояние оставлялось только если его удавалось восстановить по
доступным признакам с качеством не ниже $R^2 = 0{,}70$. По итогам этого отбора в
финальную схему вошли четыре состояния: социальная поддержка, технологический
прогресс, кредитный оптимизм и продовольственная безопасность.

\subsection{Восстановление состояний}

Для восстановления состояний использовалась регуляризованная линейная модель
Ridge:
\begin{equation}
    \hat S_j = \beta_0 + X_{cause}\beta + D_{sphere}\gamma,
\end{equation}
где $X_{cause}$ --- причины после VIF-фильтра, а $D_{sphere}$ --- dummy-признаки
сферы занятости. Регуляризация уменьшает чувствительность коэффициентов к
остаточной коррелированности признаков и к шуму в отдельных группах.

Качество восстановления оценивалось по out-of-fold-прогнозам. Это важно:
классификатор на втором этапе обучался не на состояниях, полученных моделью,
видевшей ту же строку при обучении, а на OOF-предсказаниях состояний.

\subsection{Проверка нелинейностей}

Для причин проверялся банк преобразований:
\begin{center}
\code{z, z2, z3, z4, log1p, sqrt, cbrt, log\_shift, sqrt\_shift, inv\_shift}
\end{center}
и пороговые hinge-функции по квантилям. Преобразования сравнивались с базовой
линейной моделью по изменению $R^2$. В итоговой схеме отобранные преобразования
использовались при восстановлении состояний. При этом отдельные одиночные
преобразования давали небольшой прирост, а основной практический эффект банка
проявился для наиболее сложного состояния --- индекса кредитного оптимизма.

\subsection{Сообщества}

Общественные причины внутри одного сообщества постоянны. Их комбинация задает
категориальный идентификатор \code{community\_id}, который полезен для анализа
групповой структуры данных. В ноутбуке проверялись варианты с включением
\code{community\_id} как fixed effect, однако на текущем rowwise-разбиении
лучший воспроизводимый результат дала модель со сферой занятости, но без
\code{community one-hot}. Поэтому сообщества использовались в отчете прежде
всего для интерпретации и диагностики, а не как часть финального классификатора.

\subsection{Классификатор и метрика}

На втором этапе обучалась мультиномиальная логистическая регрессия:
\begin{equation}
    P(y=k|z) = \frac{\exp(w_k^T z)}{\sum_m \exp(w_m^T z)},
\end{equation}
где $z$ включает OOF-предсказания выбранных состояний, причины после VIF-фильтра
и сферу занятости. Для регуляризации использовалась L2-норма, параметр
логистической регрессии в финальном запуске равен $C=8{,}0$.

Основная метрика --- micro precision. В однометочной многоклассовой задаче она
совпадает с accuracy:
\begin{equation}
    Precision_{micro} = \frac{\sum_k TP_k}{\sum_k (TP_k + FP_k)}.
\end{equation}

\newpage
